\chapter{Пешеходные маршруты}

Шесть маршрутов, охватывающих основные ландшафтные типы острова:
от~культовых вершин кальдеры до~морских дюн. Для всех троп
рекомендуется обувь с~жёсткой подошвой, 2~л воды летом и~головной
убор; связь спасателей --- \textbf{112}.

% --------------------------------------------------------------------
\hike{GC-1}{Рокэ-Нубло \sightt{}}{3,5~км (круг)}{220~м}{1\,ч~30\,мин}{Лёгкая}
\label{hike:gc1}

\marginfact{Старт: La~Goleta\\27.9625\degree N\\15.6372\degree W\\Высота 1580~м}

\noindent Самая посещаемая тропа Гран-Канарии --- и~по праву.
От~парковки \textit{La~Goleta} на~дороге GC-600 маркированная
песчано-каменистая тропа поднимается через сосны к~плато
\textit{Roque Nublo}. Последние 300~м --- по~открытому камню;
склон пологий, но~в~ветер --- будьте внимательны.

% --- GC-1 Rocque Nublo trail schematic + elevation profile -----------
\begin{figure}[h]
\centering
\begin{tikzpicture}[scale=0.9, every node/.style={font=\footnotesize}]
  % background route line
  \draw[baedsepia, line width=1.2pt, rounded corners]
    (0,0) -- (1.4,0.4) -- (2.6,0.2) -- (3.8,1.1) -- (4.6,1.6)
    -- (4.8,1.8) -- (4.6,1.6) -- (3.8,1.1) -- (2.6,0.2) -- (1.4,0.4) -- (0,0);
  % waypoints
  \foreach \x/\y/\n in {0/0/{1},1.4/0.4/{2},2.6/0.2/{3},3.8/1.1/{4},4.8/1.8/{5}}
    \node[circle, draw=baedred, fill=white, inner sep=1.2pt] at (\x,\y) {\scriptsize\n};
  \node[anchor=west]      at (0.1,-0.25) {La Goleta};
  \node[anchor=south]     at (4.8,1.9)   {\textbf{Роке-Нубло} 1813\,м};
  \node[anchor=west]      at (3.9,1.1)   {плато};
  \draw[->, baedmute]     (-0.6,-0.3) -- (-0.1,-0.3) node[midway, below]{\scriptsize GC-600};
\end{tikzpicture}
\hspace{0.6em}
\begin{tikzpicture}
\begin{axis}[
  width=4.8cm, height=3.0cm,
  xlabel={\scriptsize км}, ylabel={\scriptsize м},
  xmin=0, xmax=3.5, ymin=1560, ymax=1830,
  tick label style={font=\tiny},
  label style={font=\tiny},
  every axis plot/.append style={thick, baedred}]
\addplot coordinates {(0,1580) (0.7,1620) (1.6,1700) (1.9,1813) (2.2,1700) (3.5,1580)};
\end{axis}
\end{tikzpicture}
\caption*{\footnotesize Маршрут \textbf{GC-1}: схема и~профиль высот.}
\end{figure}


\waypoint{0,0~км}{Парковка La~Goleta (1580~м). Щит Senda S-70.}
\waypoint{0,7~км}{Сосновый лес; слева --- вид на~Рокэ-Бентайга.}
\waypoint{1,6~км}{Плато; справа \textit{El~Fraile} («Монах»).}
\waypoint{1,9~км}{Основание \sightt{Рокэ-Нубло} (1813~м). Разворот.}
\waypoint{3,5~км}{Возвращение в~La~Goleta.}

\begin{detour}
Расширенный вариант: спуск на~запад к~\textit{Ayacata} (3~ч, --600~м);
автобус №18 обратно в~Техеду.
\end{detour}

% --------------------------------------------------------------------
\hike{GC-2}{Крус-де-Техеда \,$\rightarrow$\, Артенара}{8,4~км}{+120\,/\,--480~м}{3\,ч}{Средняя}

\marginfact{Старт:\\Cruz de~Tejeda\\(1580~м)\\Финиш:\\Artenara\\(1270~м)}

\noindent Классический \textit{камино реаль}: древняя мощёная
дорога, по~которой крестьяне ещё в~XVIII~в. гоняли коз через
центральную кальдеру. Тропа идёт по~северному склону,
постоянно открывая вид на~\textit{Roque Bentayga}.

% --- GC-2 Tejeda -> Artenara -----------------------------------------
\begin{figure}[h]
\centering
\begin{tikzpicture}[scale=0.8, every node/.style={font=\footnotesize}]
  \draw[baedsepia, line width=1.2pt]
    (0,0) -- (1.2,0.3) -- (2.3,-0.1) -- (3.6,-0.4) -- (4.8,-0.8) -- (6.0,-1.1);
  \foreach \x/\y/\n in {0/0/{1},1.2/0.3/{2},2.3/-0.1/{3},3.6/-0.4/{4},6.0/-1.1/{5}}
    \node[circle, draw=baedred, fill=white, inner sep=1.2pt] at (\x,\y) {\scriptsize\n};
  \node[anchor=south]  at (0,0.15)    {\textbf{Cruz de Tejeda}};
  \node[anchor=north]  at (6.0,-1.2)  {\textbf{Artenara}};
  \node[anchor=west]   at (2.4,-0.15) {Degollada};
  \node[anchor=west]   at (3.7,-0.45) {Pinar de Chirimique};
\end{tikzpicture}
\hspace{0.4em}
\begin{tikzpicture}
\begin{axis}[
  width=4.6cm, height=3cm,
  xlabel={\scriptsize км}, ylabel={\scriptsize м},
  xmin=0, xmax=8.5, ymin=1200, ymax=1650,
  tick label style={font=\tiny},
  label style={font=\tiny},
  every axis plot/.append style={thick, baedred}]
\addplot coordinates {(0,1580) (2.2,1600) (4.0,1480) (6.3,1350) (8.4,1270)};
\end{axis}
\end{tikzpicture}
\caption*{\footnotesize Маршрут \textbf{GC-2}: схема и~профиль высот.}
\end{figure}


\waypoint{0,0~км}{Крус-де-Техеда, памятный крест (1580~м).}
\waypoint{2,2~км}{Смотровая \textit{Degollada de~las~Palomas}.}
\waypoint{4,0~км}{Сосновый лес \textit{Pinar de~Chirimique}.}
\waypoint{6,3~км}{Часовня \textit{Ermita de~la~Cuevita} ---
                   пещерный храм, вырубленный в~туфе.}
\waypoint{8,4~км}{Артенара, площадь с~рестораном \textit{Mirador
                   de~la~Atalaya}.}

Возврат: автобус №220 в~Лас-Пальмас (ежечасно) или №305 в~Техеду.

% --------------------------------------------------------------------
\hike{GC-3}{Тамадаба: сосновый круг \sight{}}{9,0~км (круг)}{350~м}{3\,ч~30\,мин}{Средняя}

\marginfact{Старт:\\Llanos de~la~Mimbre\\27.9902\degree N\\15.7133\degree W\\Высота 1300~м}

\noindent Самый нетронутый сосновый массив острова: деревья
возрастом 400+\,лет, густой туман по~утрам, полная тишина.
Тропа --- петля вокруг вершины \textit{Pico de~Tamadaba}
(1444~м) с~несколькими смотровыми над~западной «стеной»
острова.

% --- GC-3 Tamadaba loop ----------------------------------------------
\begin{figure}[h]
\centering
\begin{tikzpicture}[scale=0.85, every node/.style={font=\footnotesize}]
  \draw[baedsepia, line width=1.2pt, rounded corners]
    (0,0) -- (-1.0,0.4) -- (-1.4,1.4) -- (-0.6,2.2) -- (0.8,2.3)
    -- (1.9,1.7) -- (2.1,0.8) -- (1.3,0.2) -- (0,0);
  \foreach \x/\y/\n in {0/0/{1},-1.4/1.4/{2},0.8/2.3/{3},2.1/0.8/{4},1.3/0.2/{5}}
    \node[circle, draw=baedred, fill=white, inner sep=1.2pt] at (\x,\y) {\scriptsize\n};
  \node[anchor=north] at (0,-0.1)   {\textbf{Llanos de la Mimbre}};
  \node[anchor=east]  at (-1.5,1.4) {Mirador del Balcón};
  \node[anchor=south] at (0.8,2.4)  {Pico \textbf{1444\,м}};
  % Ocean to the west
  \draw[baedrule, decoration={coil, segment length=6pt, amplitude=1pt},
        decorate] (-2.6,0.5) -- (-2.6,2.0);
  \node[anchor=east, baedmute] at (-2.6,1.25) {\scriptsize океан};
\end{tikzpicture}
\hspace{0.2em}
\begin{tikzpicture}
\begin{axis}[
  width=4.6cm, height=3cm,
  xlabel={\scriptsize км}, ylabel={\scriptsize м},
  xmin=0, xmax=9.5, ymin=1250, ymax=1500,
  tick label style={font=\tiny},
  label style={font=\tiny},
  every axis plot/.append style={thick, baedred}]
\addplot coordinates {(0,1300) (2.6,1380) (5.1,1444) (7.4,1340) (9.0,1300)};
\end{axis}
\end{tikzpicture}
\caption*{\footnotesize Маршрут \textbf{GC-3}: кольцевая схема и~профиль.}
\end{figure}


\waypoint{0,0~км}{Лагерь \textit{Llanos de~la~Mimbre}.}
\waypoint{2,6~км}{Смотровая \sight{Mirador del~Balcón}:
                   вертикальный обрыв к~океану, 600~м.}
\waypoint{5,1~км}{\textit{Pico de~Tamadaba} (1444~м).}
\waypoint{7,4~км}{Форест-роуд обратно к~лагерю.}
\waypoint{9,0~км}{Финиш.}

% --------------------------------------------------------------------
\hike{GC-4}{Барранко-де-Гуаядеке}{6,0~км (туда-обратно)}{150~м}{2\,ч}{Лёгкая}

\marginfact{Старт: Cueva\\Bermeja, конец\\дороги GC-103\\(680~м)}

\noindent Прогулка по~дну ущелья --- ботаническому резервату.
Тропа ведёт мимо пещерных домов, садов инжира и~миндаля; в~феврале
цветёт розовым эндемичный тахинасте (\textit{Echium decaisnei}).

% --- GC-4 Barranco de Guayadeque ------------------------------------
\begin{figure}[h]
\centering
\begin{tikzpicture}[scale=0.85, every node/.style={font=\footnotesize}]
  % canyon walls
  \draw[baedrule, line width=0.5pt] (-0.3,0.7) -- (6.5,1.6);
  \draw[baedrule, line width=0.5pt] (-0.3,-0.7) -- (6.5,-1.6);
  % trail
  \draw[baedsepia, line width=1.2pt] (0,0) -- (1.5,0.05) -- (3.0,-0.05) -- (4.5,0.1) -- (6.0,0);
  \foreach \x/\y/\n in {0/0/{1},1.5/0.05/{2},3.0/-0.05/{3},6.0/0/{4}}
    \node[circle, draw=baedred, fill=white, inner sep=1.2pt] at (\x,\y) {\scriptsize\n};
  \node[anchor=east]  at (-0.1,0)  {\textbf{Cueva Bermeja}};
  \node[anchor=south] at (1.5,0.1) {Tagoror};
  \node[anchor=west]  at (6.1,0)   {конец тропы};
\end{tikzpicture}
\hspace{0.2em}
\begin{tikzpicture}
\begin{axis}[
  width=4.4cm, height=2.8cm,
  xlabel={\scriptsize км}, ylabel={\scriptsize м},
  xmin=0, xmax=6.5, ymin=650, ymax=850,
  tick label style={font=\tiny},
  label style={font=\tiny},
  every axis plot/.append style={thick, baedred}]
\addplot coordinates {(0,680) (1.5,720) (3.0,830) (4.5,760) (6.0,680)};
\end{axis}
\end{tikzpicture}
\caption*{\footnotesize Маршрут \textbf{GC-4}: схема ущелья и~профиль.}
\end{figure}


\waypoint{0,0~км}{Деревня \textit{Cueva Bermeja}, пещерная
                   церковь Сан-Бартоломе.}
\waypoint{1,5~км}{Ресторан \textit{Tagoror}, также вырубленный
                   в~скале; хорошая \textit{carne cabra}.}
\waypoint{3,0~км}{Конец тропы, полноводный ручей
                   (зимой); разворот.}

\begin{detour}
Ответвление на~север (1~ч дополнительно) --- к~археологическому
сайту \textit{Cuevas de~los~Canarios}, могильник гуанчей.
\end{detour}

% --------------------------------------------------------------------
\hike{GC-5}{Агаэте \,$\rightarrow$\, Тамадаба (GR-131)}{14,5~км}{+1150~м}{6\,ч}{Тяжёлая}

\marginfact{Старт: площадь\\Agaete (40~м)\\Финиш:\\Llanos de\\la~Mimbre (1300~м)}

\noindent Северный отрезок дальнего маршрута \textit{GR-131 El~Bastón}:
непрерывный подъём из~приморского городка в~высокогорный
сосновый лес. Маршрут в~один конец; возврат --- автобусом №101
через Сан-Николас.

\waypoint{0,0~км}{Агаэте, церковь \textit{Nuestra Señora de~la~Concepción}.}
\waypoint{3,2~км}{\textit{Los~Berrazales}, последние кофейные
                   плантации Европы.}
\waypoint{7,0~км}{\textit{El~Hornillo}, перевал (820~м).}
\waypoint{11,8~км}{Ребро \textit{Tirma}: первые сосны.}
\waypoint{14,5~км}{Лагерь \textit{Llanos de~la~Mimbre}.}

% --------------------------------------------------------------------
\hike{GC-6}{Дюны Маспаломаса \sightt{}}{6,3~км}{плоско}{1\,ч~45\,мин}{Лёгкая}

\marginfact{Лучше всего\\за~час до\\заката}

\noindent Прогулка по~единственному в~Европе активному дюнному
полю. Тропа не~маркирована --- ориентиры маяк и~гряды песка.
Воды --- не~менее 1\,л в~жаркий день; в~полдень песок разогревается
до~60\degree.

% --- GC-6 Dunes of Maspalomas ---------------------------------------
\begin{figure}[h]
\centering
\begin{tikzpicture}[scale=0.85, every node/.style={font=\footnotesize}]
  % ocean
  \fill[baedrule!30] (-0.3,-1.2) rectangle (7,-0.6);
  \node[baedmute] at (3.3,-0.9) {\scriptsize океан};
  % dunes (wavy)
  \draw[baedsepia, thick, decorate, decoration={coil, aspect=0, segment length=6pt, amplitude=1pt}]
       (-0.2,0) -- (6.8,0);
  % trail
  \draw[baedred, line width=1.0pt]
    (0,-0.4) -- (1.2,-0.2) -- (3.0,0.3) -- (4.6,0.2) -- (6.3,-0.4);
  \foreach \x/\y/\n in {0/-0.4/{1},1.2/-0.2/{2},3.0/0.3/{3},6.3/-0.4/{4}}
    \node[circle, draw=baedred, fill=white, inner sep=1.2pt] at (\x,\y) {\scriptsize\n};
  \node[anchor=east]  at (-0.1,-0.4) {\textbf{Faro}};
  \node[anchor=south] at (1.2,-0.15) {Charca};
  \node[anchor=south] at (3.0,0.35)  {высшая гряда};
  \node[anchor=west]  at (6.4,-0.4)  {\textbf{Playa del Inglés}};
\end{tikzpicture}
\caption*{\footnotesize Маршрут \textbf{GC-6}: песчаное побережье.}
\end{figure}


\waypoint{0,0~км}{Маяк \textit{Faro de~Maspalomas} (XIX~в., 60~м).}
\waypoint{1,2~км}{\textit{Charca de~Maspalomas}, солоноватая
                   лагуна, остановка перелётных птиц.}
\waypoint{3,0~км}{Центр дюн, высшая гряда (~12~м).}
\waypoint{6,3~км}{\textit{Playa del Inglés}, променад.}

\begin{detour}
Возврат трамвайчиком \textit{Tren Turístico} или пешком берегом ---
ещё 5~км по~мокрому песку.
\end{detour}
